\documentclass[11pt]{article}
\usepackage[margin=1in]{geometry}
\usepackage{hyperref}
\usepackage{enumitem}
\usepackage{booktabs}
\usepackage{xcolor}
\title{ProbOS --- Month 3, Week 10\\GPU Kernel Path}
\author{Nisong Monyimba}
\date{\today}

\begin{document}
\maketitle

\section*{Week 10 goal}
Build and honestly benchmark a GPU-accelerated path for
\texttt{MonteCarloEngine}'s forward-simulation loop, per Month 3's
own stated priority order and an explicit decision to finish
kernel/speed/core work before further feature work (PDSL v0.2, etc.).

\textbf{Hardware confirmed before writing this plan:} NVIDIA GeForce
RTX 3050 (4GB VRAM), CUDA 13.3 drivers already working correctly
through WSL2 via \texttt{nvidia-smi} -- no additional driver setup
needed. Genuinely sufficient for ProbOS's typical particle counts
(N = 5{,}000--100{,}000).

\textbf{Critical constraint, stated up front:} GitHub Actions CI
runners do NOT have GPU access. Any GPU-specific code must be written
to gracefully skip (via \texttt{pytest.mark.skipif}) when no GPU is
present, and CI will only validate that the GPU-path code exists and
does not break anything else -- genuine GPU correctness and
performance validation can only happen on this local machine, not CI.
This must not be papered over or silently assumed to work.

\section*{Day-by-day plan}

\subsection*{Day 1 -- Environment setup and feasibility check}
\begin{itemize}[leftmargin=*]
  \item Install CuPy matched to the confirmed CUDA 13.x toolkit
    version (CuPy chosen over raw Numba CUDA kernels: it is a
    near-drop-in NumPy replacement, meaning \texttt{BatteryModel2Cell.
    forward\_batch()}'s existing vectorised NumPy code needs minimal
    rework, consistent with keeping Week 4's original model
    unchanged wherever possible)
  \item Verify basic GPU array operations work correctly from within
    the existing \texttt{.venv} (no engine changes yet) -- a genuine
    feasibility smoke test, not assumed
  \item Confirm \texttt{cupy.cuda.is\_available()} and basic
    array-transfer round-trip correctness (CPU array -> GPU -> CPU,
    values unchanged)
\end{itemize}

\subsection*{Day 2 -- GPU-compatible forward\_batch}
\begin{itemize}[leftmargin=*]
  \item Determine whether \texttt{BatteryModel2Cell.forward\_batch()}'s
    existing NumPy operations are already CuPy-compatible as written,
    or need a thin dispatch layer (array-module parameter, following
    the common \texttt{xp = cupy if on\_gpu else numpy} pattern)
  \item Initial correctness check: same inputs, CPU (NumPy) vs GPU
    (CuPy) outputs, at an appropriate tolerance (GPU computations
    commonly default to float32 unless explicitly configured for
    float64 -- this must be checked explicitly, not assumed to match
    CPU float64 precision)
\end{itemize}

\subsection*{Day 3 -- GPUMonteCarloEngine}
\begin{itemize}[leftmargin=*]
  \item Build a new class following \texttt{MonteCarloEngine}'s exact
    existing architecture (same constructor signature, same
    \texttt{MCResult} return type extended if needed), but running the
    simulation loop with CuPy arrays on the GPU
  \item Existing \texttt{Model}/\texttt{Distribution} ABCs reused
    unchanged -- domain-agnostic kernel design already established
    should extend to the GPU path without modification
\end{itemize}

\subsection*{Day 4 -- Correctness validation}
\begin{itemize}[leftmargin=*]
  \item Full \texttt{GPUMonteCarloEngine} run vs
    \texttt{MonteCarloEngine} (CPU) run, same seed, same N, checked at
    an honestly-determined tolerance (not blindly \texttt{rtol=1e-10}
    if float32 GPU precision cannot support it -- this must be
    determined empirically, not assumed)
\end{itemize}

\subsection*{Day 5 -- Honest benchmarking}
\begin{itemize}[leftmargin=*]
  \item Measure real wall-clock speedup (or lack thereof) vs both the
    Python engine and the existing C++/OpenMP engine, across a range
    of N (small N may show GPU overhead DOMINATING actual compute --
    kernel launch and host-device transfer overhead can make small
    problems slower on GPU, not faster; this must be measured, not
    assumed away)
  \item Per the project's established ``measure, don't assume''
    discipline (Month 1 Week 4's DP profiling found one genuine
    non-win alongside two large wins) -- report the real numbers
    honestly, including if GPU turns out not to help at ProbOS's
    typical N for this specific workload
\end{itemize}

\subsection*{Day 6 -- Test coverage, pre\_week\_audit.sh, CI}
\begin{itemize}[leftmargin=*]
  \item Tests must skip gracefully in CI (no GPU) via
    \texttt{pytest.mark.skipif(not cupy\_available, ...)}, while still
    running for real locally
  \item Confirm \texttt{pre\_week\_audit.sh} and CI both pass with the
    new code present but GPU tests skipped -- CI must not silently
    report success without actually running anything meaningful for
    this feature; the skip must be visible in the test output, not
    hidden
\end{itemize}

\subsection*{Day 7 -- Documentation + retrospective}
\begin{itemize}[leftmargin=*]
  \item Update \texttt{docs/sbir\_phase1\_onepager.md}: move the GPU
    kernel item from the SBIR Ask section to demonstrated results IF
    genuinely complete and shown to help; if the honest benchmark
    result is negative or mixed, document that accurately instead of
    overselling
  \item Retrospective appended to this document once complete
\end{itemize}

\section*{Exit criteria}
A working \texttt{GPUMonteCarloEngine} exists, is proven numerically
correct against the CPU engine at an honestly-determined tolerance,
and its real performance characteristics (whether positive, negative,
or mixed across different N) are measured and reported truthfully --
not assumed, and not silently downplayed if unfavourable.

\end{document}
